% Options for packages loaded elsewhere
\PassOptionsToPackage{unicode}{hyperref}
\PassOptionsToPackage{hyphens}{url}
%
\documentclass[
]{book}
\usepackage{amsmath,amssymb}
\usepackage{lmodern}
\usepackage{iftex}
\ifPDFTeX
  \usepackage[T1]{fontenc}
  \usepackage[utf8]{inputenc}
  \usepackage{textcomp} % provide euro and other symbols
\else % if luatex or xetex
  \usepackage{unicode-math}
  \defaultfontfeatures{Scale=MatchLowercase}
  \defaultfontfeatures[\rmfamily]{Ligatures=TeX,Scale=1}
\fi
% Use upquote if available, for straight quotes in verbatim environments
\IfFileExists{upquote.sty}{\usepackage{upquote}}{}
\IfFileExists{microtype.sty}{% use microtype if available
  \usepackage[]{microtype}
  \UseMicrotypeSet[protrusion]{basicmath} % disable protrusion for tt fonts
}{}
\makeatletter
\@ifundefined{KOMAClassName}{% if non-KOMA class
  \IfFileExists{parskip.sty}{%
    \usepackage{parskip}
  }{% else
    \setlength{\parindent}{0pt}
    \setlength{\parskip}{6pt plus 2pt minus 1pt}}
}{% if KOMA class
  \KOMAoptions{parskip=half}}
\makeatother
\usepackage{xcolor}
\usepackage{longtable,booktabs,array}
\usepackage{calc} % for calculating minipage widths
% Correct order of tables after \paragraph or \subparagraph
\usepackage{etoolbox}
\makeatletter
\patchcmd\longtable{\par}{\if@noskipsec\mbox{}\fi\par}{}{}
\makeatother
% Allow footnotes in longtable head/foot
\IfFileExists{footnotehyper.sty}{\usepackage{footnotehyper}}{\usepackage{footnote}}
\makesavenoteenv{longtable}
\usepackage{graphicx}
\makeatletter
\def\maxwidth{\ifdim\Gin@nat@width>\linewidth\linewidth\else\Gin@nat@width\fi}
\def\maxheight{\ifdim\Gin@nat@height>\textheight\textheight\else\Gin@nat@height\fi}
\makeatother
% Scale images if necessary, so that they will not overflow the page
% margins by default, and it is still possible to overwrite the defaults
% using explicit options in \includegraphics[width, height, ...]{}
\setkeys{Gin}{width=\maxwidth,height=\maxheight,keepaspectratio}
% Set default figure placement to htbp
\makeatletter
\def\fps@figure{htbp}
\makeatother
\setlength{\emergencystretch}{3em} % prevent overfull lines
\providecommand{\tightlist}{%
  \setlength{\itemsep}{0pt}\setlength{\parskip}{0pt}}
\setcounter{secnumdepth}{5}
\usepackage{booktabs}
\ifLuaTeX
  \usepackage{selnolig}  % disable illegal ligatures
\fi
\usepackage[]{natbib}
\bibliographystyle{apalike}
\IfFileExists{bookmark.sty}{\usepackage{bookmark}}{\usepackage{hyperref}}
\IfFileExists{xurl.sty}{\usepackage{xurl}}{} % add URL line breaks if available
\urlstyle{same} % disable monospaced font for URLs
\hypersetup{
  pdftitle={Biology for environmental management pocketguide},
  pdfauthor={cjlortie},
  hidelinks,
  pdfcreator={LaTeX via pandoc}}

\title{Biology for environmental management pocketguide}
\author{cjlortie}
\date{}

\begin{document}
\maketitle

{
\setcounter{tocdepth}{1}
\tableofcontents
}
\hypertarget{introduction}{%
\chapter{Introduction}\label{introduction}}

\includegraphics[width=4in,height=\textheight]{./bolt.png}

Welcome to biology for environmental management. The goal of this course is to examine global grand challenges that we face through the lens of solutions from science. Science can be magical because it provides us with incredible outcomes in so many domains of the human endeavour. Here, we work to find the magic spells from science for the environment to promote the well-being of people and the planet.

There is an opportunity for us to mobilize science to do social good. We we can leverage what we know, what we do not know, i.e.~uncertainty, and what we need to know, to enable conservation, protection, and better management of the environment. This benefits the environment. It benefits us. Here are the \href{https://www.facetsjournal.com/doi/10.1139/facets-2020-0021}{ten simple rules} the underpin an evidence-informed philosophy for the entire course \citep{RN6861}.

\hypertarget{course-outline}{%
\subsection*{Course outline}\label{course-outline}}
\addcontentsline{toc}{subsection}{Course outline}

If you are electing to engage with this learning opportunity formally, please see the official course outline for specific details.

There are three summative assessments.\\
1. Test (apply the challenge-solution framework).\\
2. Make a scientific comic or infographic for a new challenge of your choice.\\
3. Write a super short ignite synthesis paper.

\hypertarget{learning-outcomes}{%
\subsection*{Learning outcomes}\label{learning-outcomes}}
\addcontentsline{toc}{subsection}{Learning outcomes}

\begin{enumerate}
\def\labelenumi{\arabic{enumi}.}
\tightlist
\item
  Critically read environmental science peer-reviewed journal publications.\\
\item
  Reverse-engineer the critical reproducible science tools using peer-reviewed publications.\\
\item
  Appreciate the extent and scope of environmental challenges we face globally.\\
\item
  Explain the balance between direct human needs and environmental health.\\
\item
  Do a formal synthesis such as meta-analysis or systematic review.\\
\item
  Effectively communicate scientific synthesis findings to the public.
\end{enumerate}

\hypertarget{steps}{%
\subsection*{Steps}\label{steps}}
\addcontentsline{toc}{subsection}{Steps}

\hypertarget{module-1.}{%
\subsubsection*{Module 1.}\label{module-1.}}
\addcontentsline{toc}{subsubsection}{Module 1.}

Read a total of 9 useful peer-reviewed science publications.\\
Test your practical knowledge by applying to a new challenge.

\hypertarget{module-2.}{%
\subsubsection*{Module 2.}\label{module-2.}}
\addcontentsline{toc}{subsubsection}{Module 2.}

Choose your own adventure (i.e.~a dimension of an environmental challenge you care about).\\
Draw a comic or infographic to communicate challenge to the public.\\
Write a short synthesis paper on this topic for a scientific audience.

\hypertarget{rationale}{%
\subsubsection*{Rationale}\label{rationale}}
\addcontentsline{toc}{subsubsection}{Rationale}

For each environmental management challenge case examined, students will be responsible for reading the literature provided at their own pace. The professor will facilitate learning as needed.

The goal is to become more literate environmental citizens and develop, consolidate, and evaluate critical environmental science thinking and problem solving.

The first module highlights some of the most pressing challenges and more common replicable tools used by the scientific community. The summative test is provided immediately at the start of course to enable asynchronous work and provide a clear, transparent, and testable outcome for this module.

The second module provides an opportunity for students in this upper-year offering to do a deep dive into a topic that care about deeply. The dimension of the challenge and the solution they pick is open provided it is well articulated. The graphical assignment is a stepping stone or scaffolding to the final paper. It is also a chance to be as creative as students elect to be with communicating science to the public. The final paper is an Ignite, Forum, or Mini-review format contribution on their topic appropriate for a general science journal. These types of papers are increasingly common and important in science and used extensively for evidence-informed decision making by leaders.

\hypertarget{citation}{%
\subsection*{Citation}\label{citation}}
\addcontentsline{toc}{subsection}{Citation}

Lortie, CJ (2021): Biology for environmental management pocketguide. figshare. Book. \url{https://doi.org/10.6084/m9.figshare.15031752.v4}

\hypertarget{license}{%
\subsection*{License}\label{license}}
\addcontentsline{toc}{subsection}{License}

This work is licensed under a Creative Commons Attribution-NonCommercial-ShareAlike 4.0 International License.

\hypertarget{topics}{%
\subsection*{Topics}\label{topics}}
\addcontentsline{toc}{subsection}{Topics}

\includegraphics[width=3in,height=\textheight]{./outline.png}

Here is an overview of the content and topics covered in this course of study. Complete are your own pace, asynchronously. However, please check the official course outline if you are doing the work for credit to ensure you submit summative work at the appropriate times.

\hypertarget{instructions}{%
\subsection*{Instructions}\label{instructions}}
\addcontentsline{toc}{subsection}{Instructions}

Read and and use the papers to link environmental challenges that we collectively face with potential solutions. Only one solution per challenge is suggested here, but there are many dimensions to each challenge and numerous solutions too.

The link to decks are optional. They are my interpretation of the papers from a science-to-magic philosophy and identify the salient elements and concepts from each reading that resonated with my perspective as an ecologist.

\hypertarget{schedule}{%
\subsection*{Schedule}\label{schedule}}
\addcontentsline{toc}{subsection}{Schedule}

\begin{tabular}{rll}
\toprule
week & grand challenges & solutions\\
\midrule
1 & {}[Ten simple rules to facilitate evidence implementation in the environmental sciences](https://www.facetsjournal.com/doi/10.1139/facets-2020-0021) & this course :)\\
2 & {}[climate change](https://agupubs.onlinelibrary.wiley.com/doi/abs/10.1029/2010GL042845) \& [deck](https://figshare.com/articles/presentation/BIOL4265\_climate\_change/14960109) & {}[R](https://ojs.library.queensu.ca/index.php/IEE/article/view/6559) \& [deck](https://figshare.com/articles/presentation/BIOL4265\_R\_is\_a\_solution\_/14960121)\\
3 & {}[pollinator declines](http://science.sciencemag.org/content/339/6127/1611) \& [deck](https://figshare.com/articles/presentation/BIOL4265\_pollinator\_declines/14960130) & {}[citizen science](https://www.sciencedirect.com/science/article/pii/S0006320716301963?via\%3Dihub) \& [deck](https://figshare.com/articles/presentation/BIOL4265\_citizen\_science/14960139)\\
4 & {}[invasive species](https://link.springer.com/article/10.1007/s10530-017-1461-x) \& [deck](https://figshare.com/articles/presentation/BIOL4265\_invasive\_species/14960145) & {}[team science](https://esajournals.onlinelibrary.wiley.com/doi/full/10.1890/120370) \& [deck](https://figshare.com/articles/presentation/BIOL4265\_team\_science/14960154)\\
5 & {}[land use challenges](https://esajournals.onlinelibrary.wiley.com/doi/full/10.1002/ecs2.2367) \& [deck](https://figshare.com/articles/presentation/BIOL4265\_landuse\_challenges/14960337) & {}[systematic review](https://onlinelibrary.wiley.com/doi/abs/10.1111/j.1600-0706.2013.00970.x) \& [deck](https://figshare.com/articles/presentation/BIOL4265\_systematic\_reviews/14962560)\\
\addlinespace
6 & prep for test & {}[ten simple rules for decision making](https://journals.plos.org/ploscompbiol/article?id=10.1371/journal.pcbi.1007706)\\
7 & {}[take-home test](https://www.facetsjournal.com/doi/full/10.1139/facets-2020-0108) & select a dimension of challenge, find a replicable solution, use that paper here\\
8 & choose your own adventure (i.e. challenge): search lit, read, sketch ideas & {}[synthesis tools in science](https://onlinelibrary.wiley.com/doi/full/10.1111/j.1600-0706.2013.00970.x)\\
9 & complete \& submit scientific comic or infographic & {}[scientific communication](https://journals.plos.org/ploscompbiol/article?id=10.1371/journal.pcbi.1005845)\\
10 & choose your own adventure: write paper (review, systematic review, or meta) & {}[meta-analysis](https://onlinelibrary.wiley.com/doi/abs/10.1002/jrsm.1109)\\
\addlinespace
11 & choose your own adventure: write paper & {}[useful rules for science papers](https://journals.plos.org/ploscompbiol/article?id=10.1371/journal.pcbi.1005619)\\
12 & choose your own adventure: submit final synthesis paper & {}[Guidelines for authors in a leading journal in short-format papers](https://www.oikosjournal.org/authors/author-guidelines) \& [deck](https://figshare.com/articles/presentation/A\_summary\_of\_short\_scientific\_synthesis\_contributions/17069735)\\
\bottomrule
\end{tabular}

\hypertarget{climate}{%
\chapter{Climate change}\label{climate}}

\includegraphics[width=3in,height=\textheight]{./climate.png}

\hypertarget{learning-outcomes-1}{%
\subsubsection*{Learning outcomes}\label{learning-outcomes-1}}
\addcontentsline{toc}{subsubsection}{Learning outcomes}

\begin{enumerate}
\def\labelenumi{\arabic{enumi}.}
\tightlist
\item
  Explore one set of dimensions associated with a changing climate.
\item
  Link science for cities to your life.\\
\item
  Explore one tool that can enable replicable solutions.
\end{enumerate}

\hypertarget{context}{%
\subsubsection*{Context}\label{context}}
\addcontentsline{toc}{subsubsection}{Context}

We experience weather but live with climate. Climate is complex. Climate change has both effects on us and is a response to many drivers including anthropenic processes. The reading provided examines urban effects specifically \citep{RN5995}. Use the ten simple rules suggested to structure your analysis of this paper \citep{RN6861}. There is at least one solution that can enable better science. Replication and being able to test the same ideas again can be done using an open-source, and free, programming language to work with data, draw plots, and do statistics. \href{https://www.r-project.org}{R} is one such tool, and it can be used to promote solutions for others to try with their challenges because the code can be shared (i.e.~fondly recall math classes, show your work), and this documentation of science using data improves our solutions \citep{RN4523}.

\hypertarget{reflection-questions}{%
\subsubsection*{Reflection questions}\label{reflection-questions}}
\addcontentsline{toc}{subsubsection}{Reflection questions}

\begin{enumerate}
\def\labelenumi{\arabic{enumi}.}
\tightlist
\item
  What are some of the other dimensions of climate change?
\item
  What tools or other solutions do scientists in many fields use to address a changing climate?
\item
  What is the biggest challenge of this issue that must be addressed and would R help?
\end{enumerate}

\hypertarget{formative-checklist-or-steps}{%
\subsubsection*{Formative checklist or steps}\label{formative-checklist-or-steps}}
\addcontentsline{toc}{subsubsection}{Formative checklist or steps}

\begin{enumerate}
\def\labelenumi{\arabic{enumi}.}
\tightlist
\item
  Read the paper.\\
\item
  Consider the questions provided, to guide your thinking and practice implementing evidence to do magic (i.e.~potentially use science to address challenges). These are not graded, and the purpose is to reflect and actively engage with readings.\\
\item
  Review the slide decks (optional) after you read and consider the two papers to see if similar concepts resonated with you.
\end{enumerate}

\hypertarget{pollinators}{%
\chapter{Pollinator declines}\label{pollinators}}

\includegraphics[width=3in,height=\textheight]{./bee.png}

\hypertarget{learning-outcomes-2}{%
\subsubsection*{Learning outcomes}\label{learning-outcomes-2}}
\addcontentsline{toc}{subsubsection}{Learning outcomes}

\begin{enumerate}
\def\labelenumi{\arabic{enumi}.}
\tightlist
\item
  Appreciate the relative importance of pollinators globally.
\item
  Identify the links between pollinators, plants, and functional ecology.\\
\item
  Explore one tool that can enable replicable solutions.
\end{enumerate}

\hypertarget{context-1}{%
\subsubsection*{Context}\label{context-1}}
\addcontentsline{toc}{subsubsection}{Context}

We need to eat. Pollinators provide many fundamental ecological services including interactions with plants that we use for crops and sustenance. Most flowering plants are not wind pollinated. This reading explores co-occurrence patterns between plants and pollinators to document long-term declines \citep{RN6388}. Not great news. However, there are solutions - many. One of the most exciting to me is citizen science. Citizen science is a powerful tool for conservationists \citep{RN5997}. It is also profound because we are all citizens, together, and scientific thinking comes naturally to humans. We love to count, observe natural systems, and have strong inherent tendencies to connect with nature. This is a solution for many global grand challenges associated with the environment.

\hypertarget{reflection-questions-1}{%
\subsubsection*{Reflection questions}\label{reflection-questions-1}}
\addcontentsline{toc}{subsubsection}{Reflection questions}

\begin{enumerate}
\def\labelenumi{\arabic{enumi}.}
\tightlist
\item
  Does the species/taxa of pollinators matter?\\
\item
  Citizen science can measure processes, but can it also be used to help conserve or protect or restore too?\\
\item
  How can citizen science at universities be better enabled? How do we link what we learn with what we do?
\end{enumerate}

\hypertarget{formative-checklist-or-steps-1}{%
\subsubsection*{Formative checklist or steps}\label{formative-checklist-or-steps-1}}
\addcontentsline{toc}{subsubsection}{Formative checklist or steps}

\begin{enumerate}
\def\labelenumi{\arabic{enumi}.}
\tightlist
\item
  Read the paper.\\
\item
  Consider the questions provided, to guide your thinking and practice implementing evidence to do magic (i.e.~potentially use science to address challenges). These are not graded, and the purpose is to reflect and actively engage with readings.\\
\item
  Review the slide decks (optional) after you read and consider the two papers to see if similar concepts resonated with you.
\end{enumerate}

\hypertarget{invasions}{%
\chapter{Invasive species}\label{invasions}}

\includegraphics[width=3in,height=\textheight]{./invasion.png}

\hypertarget{learning-outcomes-3}{%
\subsubsection*{Learning outcomes}\label{learning-outcomes-3}}
\addcontentsline{toc}{subsubsection}{Learning outcomes}

\begin{enumerate}
\def\labelenumi{\arabic{enumi}.}
\tightlist
\item
  Be able to define invasive species.
\item
  Describe some of the ways that ecologists test invasive species.\\
\item
  Examine whether a grand challenge needs a small or large (even global) team.
\end{enumerate}

\hypertarget{context-2}{%
\subsubsection*{Context}\label{context-2}}
\addcontentsline{toc}{subsubsection}{Context}

Invasive species are a global issue. Introduced species that become invasive can dramatically and irrevocably modify ecosystems. The paper provided here is an exceptional example of leveraging fundamental science, i.e.~biogeography, with the contemporary study of this issue \citep{RN5996}. To study a species from one place in another, particularly between continents, teams in science can do the work more efficiently and more effectively. Team science is thus a logical solution to this issue \citep{RN6000}. As you know from your science work, collaboration is not always easy though particularly when the work is being evaluated - whether in a class for groupwork or by a university of its researchers or by funding agencies in contrasting the merit of a researcher. Working with others to tackle global grand challenges is critical, but we also need to adapt how we value work in teams in science.

\hypertarget{reflection-questions-2}{%
\subsubsection*{Reflection questions}\label{reflection-questions-2}}
\addcontentsline{toc}{subsubsection}{Reflection questions}

\begin{enumerate}
\def\labelenumi{\arabic{enumi}.}
\tightlist
\item
  What is biogeography?\\
\item
  Are invasive species always `bad' or problematic?\\
\item
  How we up-value scientists that are more inclusive and collaborative?
\end{enumerate}

\hypertarget{formative-checklist-or-steps-2}{%
\subsubsection*{Formative checklist or steps}\label{formative-checklist-or-steps-2}}
\addcontentsline{toc}{subsubsection}{Formative checklist or steps}

\begin{enumerate}
\def\labelenumi{\arabic{enumi}.}
\tightlist
\item
  Read the paper.\\
\item
  Consider the questions provided, to guide your thinking and practice implementing evidence to do magic (i.e.~potentially use science to address challenges). These are not graded, and the purpose is to reflect and actively engage with readings.\\
\item
  Review the slide decks (optional) after you read and consider the two papers to see if similar concepts resonated with you.
\end{enumerate}

\hypertarget{landuse}{%
\chapter{Landuse challenges}\label{landuse}}

\includegraphics[width=3in,height=\textheight]{./landuse.png}

\hypertarget{learning-outcomes-4}{%
\subsubsection*{Learning outcomes}\label{learning-outcomes-4}}
\addcontentsline{toc}{subsubsection}{Learning outcomes}

\begin{enumerate}
\def\labelenumi{\arabic{enumi}.}
\tightlist
\item
  Be able to describe and differentiate restoration from other challenge \& solution categories.
\item
  Explore one set of ideas that link land use to people.\\
\item
  Understand some of the variety of synthesis tools available to contemporary scientists.
\end{enumerate}

\hypertarget{context-3}{%
\subsubsection*{Context}\label{context-3}}
\addcontentsline{toc}{subsubsection}{Context}

Land is a commodity. We compete with plants and other animals for space. Land literally provides the space for most components of human society through provision of direct and indirect ecosystem services. In the challenge reading, a synthesis of the restoration research for drylands in California was summarized \citep{RN5980}. The technique or replicable solution used was a formal systematic review, but there are many such synthesis tools \citep{RN3216}. Synthesis tools provide the big picture to frame questions and evaluate the relevant merit or at least frequency that scientists use one technique versus another.

\hypertarget{reflection-questions-3}{%
\subsubsection*{Reflection questions}\label{reflection-questions-3}}
\addcontentsline{toc}{subsubsection}{Reflection questions}

\begin{enumerate}
\def\labelenumi{\arabic{enumi}.}
\tightlist
\item
  Is better late than never possible for many other ecosystems in addition to drylands?\\
\item
  Why was the direct human dimension of restoration in drylands not examined in this case study challenge?\\
\item
  What synthesis technique described in the solution paper best matched your thinking style?
\end{enumerate}

\hypertarget{formative-checklist-or-steps-3}{%
\subsubsection*{Formative checklist or steps}\label{formative-checklist-or-steps-3}}
\addcontentsline{toc}{subsubsection}{Formative checklist or steps}

\begin{enumerate}
\def\labelenumi{\arabic{enumi}.}
\tightlist
\item
  Read the paper.\\
\item
  Consider the questions provided, to guide your thinking and practice implementing evidence to do magic (i.e.~potentially use science to address challenges). These are not graded, and the purpose is to reflect and actively engage with readings.\\
\item
  Review the slide decks (optional) after you read and consider the two papers to see if similar concepts resonated with you.
\end{enumerate}

\hypertarget{testprep}{%
\chapter{Test prep}\label{testprep}}

\includegraphics[width=3in,height=\textheight]{./testprep.png}

\hypertarget{learning-outcomes-5}{%
\subsubsection*{Learning outcomes}\label{learning-outcomes-5}}
\addcontentsline{toc}{subsubsection}{Learning outcomes}

\begin{enumerate}
\def\labelenumi{\arabic{enumi}.}
\tightlist
\item
  Be able to select your own solution to an environmental challenge.
\item
  Appreciate that there are many dimensions to every challenge.\\
\item
  Write a clear, short description linking a replicable tool to a challenge studied in the scientific literature.
\end{enumerate}

\hypertarget{context-4}{%
\subsubsection*{Context}\label{context-4}}
\addcontentsline{toc}{subsubsection}{Context}

Use the `ten simple rules to facilitate evidence implementation' \citep{RN6861} and `ten simple rules for more objective decision-making' \citep{RN7223} to structure notes and test prep.

\hypertarget{steps-for-test-prep}{%
\subsection*{Steps for test prep}\label{steps-for-test-prep}}
\addcontentsline{toc}{subsection}{Steps for test prep}

\begin{enumerate}
\def\labelenumi{\arabic{enumi}.}
\tightlist
\item
  Read the challenge paper provided in the test chapter herein.\\
\item
  Make a list of all the solutions that they used or could have considered.
\item
  Select one solution (a novel one not discussed yet), find 1-2 papers about that solution and 1-2 similar papers to the challenge, see test questions.\\
\item
  Re-read the \href{https://www.facetsjournal.com/doi/10.1139/facets-2020-0021}{ten simple rules} that underpins this course and read the \href{https://journals.plos.org/ploscompbiol/article?id=10.1371/journal.pcbi.1007706}{ten simple rules for making better decisions} too.
\item
  Go to next chapter in this resource for the test.
\end{enumerate}

\hypertarget{general-worflow-for-take-home-tests}{%
\subsection*{General worflow for take-home tests}\label{general-worflow-for-take-home-tests}}
\addcontentsline{toc}{subsection}{General worflow for take-home tests}

Traditional formative assessments such as in-class tests or exams are well understood. The rules of the game and extent of practice and experience with them vary from person-to-person, but most schools still adopt fixed-duration testing with limited to no access to outside materials. Open-book tests are common as well and typically include some access to resources such as notes or course materials but can also include open and online resources. The take-home test is a powerful and positive summative assessment instrument. However, effective study and prep can not always be clearly described for the learner. Strengths can include full and open access to materials online, experts, collaborative input, and adequate time to do deep work and think with reflection on insightful answers. Nonetheless, a limitation can be divergent expectations between student or learner and the graders. A take-home test is a bit of a chimera because it includes components of routine work, thinking, and depth associated with essays and papers but is nonetheless at times designed more as test with specific questions, Hence, a short primer on some ideas to consider from a workflow perspective are provided herein.

\begin{enumerate}
\def\labelenumi{\arabic{enumi}.}
\tightlist
\item
  Read the test in advance and give yourself time to ruminate, reflect, and ideate.\\
\item
  Review test questions and ensure that you read all questions to the end before starting work.\\
\item
  If there are readings, compile and collate them in advance so you have everything you need in one place.\\
\item
  Review the rubric.\\
\item
  Reverse-engineer this rubric to populate a list of terms, ideas, or visual-text hooks (keywords or concepts) that a grader will likely seek in refereeing your work.
\item
  Open book = open information. Take advantage of related ideas, papers, reviews, and commit the appropriate and most parsimonious amount of time to secure supporting resources.\\
\item
  Set a time limit for specific tasks and block your time into meaningful chunks with clear intermediate-level goals. For instance, get papers in one block, review rubric in another, etc.\\
\item
  Write, take a break, and check rubric. Check your writing for clarity. If there are synthesis, design, or principle responses, step back, and consider high-level critical thinking elements or rules associated with \href{https://journals.plos.org/ploscompbiol/article?id=10.1371/journal.pcbi.1010139}{good research planning}.\\
\item
  Consider triage answers that do not match your cognitive style, expertise, or are dissatisfied with the response. Use point value and relative weighting to decide if there better ways to allocate your time or select between questions to submit if provided.\\
\item
  Simple, but review the instructions on formatting, submission, and guidelines for the format of the work. Submit.
\end{enumerate}

\hypertarget{test}{%
\chapter{Test}\label{test}}

\includegraphics[width=3in,height=\textheight]{./test.png}

\hypertarget{steps-1}{%
\subsection*{Steps}\label{steps-1}}
\addcontentsline{toc}{subsection}{Steps}

\begin{enumerate}
\def\labelenumi{\arabic{enumi}.}
\tightlist
\item
  Read \href{https://www.facetsjournal.com/doi/full/10.1139/facets-2020-0108}{this paper} \citep{RN7311}.
\item
  Answer the following questions.
\item
  Use \textbf{no more than on page per answer} with at least an 11point font, 1inch margins, and single-spacing.
\end{enumerate}

\hypertarget{test-questions}{%
\subsection*{Test questions}\label{test-questions}}
\addcontentsline{toc}{subsection}{Test questions}

\textbf{1.} Describe the environmental management challenge case associated with this topic and publication. In the summary, please ensure that you cover the following:

\begin{enumerate}
\def\labelenumi{(\alph{enumi})}
\tightlist
\item
  link the basic science to the management (i.e.~explain what the authors tested \& measured and then how this can relate to management),\\
\item
  highlight the extent (i.e.~explain if it is a global, regional, or local issue and why), and\\
\item
  propose the implication of ignoring this issue.
\end{enumerate}

Q total = \textbf{10} points.

\textbf{2.} Describe why this is not an environmental management challenge -- i.e.~play the `devil's advocate' because we know this happens all the time or more fairly every environmental issue is a balance between direct and indirect benefits from a natural system. Do not use fallacious arguments or \emph{ad hominem} to do this. Use the following lines of argument:

\begin{enumerate}
\def\labelenumi{(\alph{enumi})}
\tightlist
\item
  list 2 critical limitations of the study (or what viable improvements are needed),\\
\item
  highlight the direct needs of people for the system/idea/issue and explain,\\
\item
  explain a `hope for the best' strategy that maximizes direct benefit to humans and propose an argument for the resilience/recovery/buffering of the system without a big change (i.e.~business as usual model).
\end{enumerate}

Q total = \textbf{10} points.

\textbf{3.} Reverse engineer a general, replicable solution to this challenge.

\begin{enumerate}
\def\labelenumi{(\alph{enumi})}
\tightlist
\item
  list solutions beginning with the most evident in descending relevance order (with your novel one listed firstly),\\
\item
  propose and explain the replicable solution that is most relevant based on research,\\
\item
  cite the research you did to design your replicable solution by providing the citation to at least two solution papers and include just a single sentence for each explaining why you selected that paper as evidence. Additionally, include 1-2 papers, cited, that show it is an important and general challenge (see rubric). Ideal - a total of 4 citations, minimum 2 (depending on how well the citations fit with how you defined the bio4enviro challenge-to-solution workflow).
\end{enumerate}

Q total = \textbf{20} points.

\hypertarget{instructions-1}{%
\subsection*{Instructions}\label{instructions-1}}
\addcontentsline{toc}{subsection}{Instructions}

\begin{enumerate}
\def\labelenumi{\arabic{enumi}.}
\tightlist
\item
  Write in full sentences, no more than 1 page per answer.
\item
  Total of 3 pages :).\\
\item
  Include your name, surname, and student ID at the top of each page.
\item
  Work in your preferred text editor, but you must submit responses as a single file in PDF format only to \href{https://www.turnitin.com}{turnitin.com}. Class ID and key provided in the official course outline.
\end{enumerate}

\hypertarget{rubric}{%
\subsection*{Rubric}\label{rubric}}
\addcontentsline{toc}{subsection}{Rubric}

\begin{tabular}{rllr}
\toprule
item & concept & description & value\\
\midrule
1 & Q1a basic science & explain what was tested and measured & 4\\
2 & Q1a basic science & explain how this study relates to environmental management & 2\\
3 & Q1b extent & explain extent - global regional local and WHY & 2\\
4 & Q1c ignorance is bliss & propose implication of ignoring issue altogether and not managing & 2\\
5 & Q2a limitations & list 2 scientific limitations of study & 4\\
\addlinespace
6 & Q2b direct needs & explain what people need from this system directly & 2\\
7 & Q2c strategy & propose a hope for best minimal management strategy for challenge & 4\\
8 & Q3a solutions & list at least 3 big picture reproducible solutions or tools and explain how each can be applied to this challenge & 6\\
9 & Q3b top solution & explain the top solution or tool to challenge based on your research i.e. at least one other paper & 10\\
10 & Q3c citations & cite at least 2 (ideally 4) recent relevant papers - 1 or even 2 solution papers and 1-2 similar challenge papers in addition to the paper provided & 4\\
\addlinespace
11 & total & sum of above & 40\\
\bottomrule
\end{tabular}

\hypertarget{synthesis}{%
\chapter{Synthesis science}\label{synthesis}}

\includegraphics[width=3in,height=\textheight]{./adventure.png}

\hypertarget{learning-outcomes-6}{%
\subsubsection*{Learning outcomes}\label{learning-outcomes-6}}
\addcontentsline{toc}{subsubsection}{Learning outcomes}

\begin{enumerate}
\def\labelenumi{\arabic{enumi}.}
\tightlist
\item
  Develop your capacity to review scientific literature and scope critical ideas.
\item
  Turn your passion and interests into formal scientific evidence for social good.\\
\item
  Synthesize peer-reviewed science papers.
\end{enumerate}

\hypertarget{context-5}{%
\subsubsection*{Context}\label{context-5}}
\addcontentsline{toc}{subsubsection}{Context}

Time to choose your own adventure. Flex your creativity. Apply your critical thinking skills to an environmental challenge that is important - to \textbf{you}! Find your science-to-magic workflow and do a scientific synthesis. The \href{https://sdgs.un.org/goals}{United Nations Sustainable Development Goals (i.e.~SDGs)} are a perfect place to begin. There are a total of 17 goals describing 169 targets as of 2021. Virtually every single goal is directly or indirectly tied to the environment. Some, are very direct such as clean water, life on land, life below water, climate action, zero hunger, and poverty. We need natural systems to sustain and support these goals. Natural systems now need us, and now you here, to summarize the scientific evidence for some focused element of a goal. This is a just a suggested focus - any environmental management challenge is fair game for synthesis and certainly will benefit from the process.

\hypertarget{synthesis-1}{%
\subsubsection*{Synthesis}\label{synthesis-1}}
\addcontentsline{toc}{subsubsection}{Synthesis}

Scientific synthesis is big picture science that describe a set of studies on a pinpointed topic \citep{RN3216}. There are at least three simple and direct options that are amenable to a capture of the research associated with the environmental challenge you identified.

\begin{enumerate}
\def\labelenumi{(\arabic{enumi})}
\item
  Narrative review as a highlight, short commentary, or new idea paper that is a snapshot of the key findings from a field research summarizing the main discoveries and/or listing the most critical research gaps. Papers like these are often called \textbf{Insights}, \textbf{Forum}, or \textbf{Ignite}.
\item
  Systematic reviews are similar to narrative reviews, but clear criteria are listed explaining how you selected papers (i.e.~these search terms used in the Web of Science, only studies that had these key inclusion attributes, etc.). Systematic are more replicable because others could follow your steps and get the same set of papers and hopefully reach similar conclusions about the corpus of evidence. These reviews also typically provide some simple quantitative data about the set of studies such as number of countries where the research was done, total sample sizes, number of variables examined, or any attributes that describe what the research was for a specific detail. The narrative component is similar to the first option because it can state what we know and what do not but these reviews do so \emph{much} more precisely. Even a few numbers go a long way to convincing people about the extent that we know or have studied a subject in science. Papers like these are often termed \textbf{Short Commentary} or \textbf{Mini-Review}.
\item
  Meta-analyses are systematic reviews plus for each primary study you summarize, you capture the relative efficacy of the treatment tested. Papers like this are often termed \textbf{Reviews} or \textbf{Meta-analyses} but other terms can be used too. Note: in some fields of research the terms `systematic review' and `meta-analysis' are used interchangeably, but in most environmental sciences, meta-analyses always have a measure of the strength of evidence from each studied included in the synthesis whilst systematic reviews typically do not.
\end{enumerate}

So for instance, narrative review might provide insights into vaccine research and reported that we have tested three vaccine types but need to test more alternatives. A systematic review would state this too but mention how they checked the science, i.e.~we checked 100 papers using these terms x,y, \& z in Web of Science, and it might also state how many people were tested in total in all studies. This is thus a more powerful synthesis. However, the gold standard would be the meta-analysis that can summarize all of the above but also reports how well each vaccine type tested actually works on average across all the studies for each one.

\hypertarget{summary-of-options}{%
\subsubsection*{Summary of options}\label{summary-of-options}}
\addcontentsline{toc}{subsubsection}{Summary of options}

\begin{tabular}{llll}
\toprule
synthesis & elements & benefits & limitations\\
\midrule
narrative review & summary, insights, next steps & more opinion, can be shorter, less detailed processing the research literature & can be less compelling without some specific evidence listed and is difficult to repeat\\
systematic review & summary, insights, next steps, explanation of how studies were selected in the review, can have counts of ideas tested & more specific, can be repeated, and can be more compelling & need to keep track of how you selected papers, need to sort through papers in more detail in addition to capturing main ideas\\
meta-analyses & all of above possible but must also include an assessment of the strength of evidence of each study included in the synthesis & gold standard, reader can get a sense of how effective a treatment or intervention is relative to another & need to extract means or measures of efficacy with sample sizes or variance from each study\\
\bottomrule
\end{tabular}

\hypertarget{steps-2}{%
\subsubsection*{Steps}\label{steps-2}}
\addcontentsline{toc}{subsubsection}{Steps}

\begin{enumerate}
\def\labelenumi{\arabic{enumi}.}
\tightlist
\item
  Review the 17 SDGs.
\item
  Familiarize yourself again with the scientific synthesis options that can describe and capture the state-of-the-art research for any environmental challenge.\\
\item
  Select a topic that you care about.\\
\item
  Check Google Scholar and/or The Web of Science trying different terms.\\
\item
  If too many studies still, filter than to most cited or the last three years of research only.\\
\item
  Get those papers.\\
\item
  Do a quick read of at least the abstracts.\\
\item
  Take a break.\\
\item
  Reflect on the environmental challenge you want to tackle and ensure the terms you used and the papers you have are studying the dimension of the challenge you want to summarize.\\
\item
  Go for it. Decide if it makes sense to do a short narrative, systematic, or meta.\\
\item
  Compile whatever evidence you need from the studies to do a,b, or c (ie ideas and concepts, count up factors or replicates, or for a meta, do they all or most report similar data and list a mean and sample size you can use).\\
\item
  Move onto next step - infographic or comic.
\end{enumerate}

\hypertarget{outcomes-from-this-work}{%
\subsubsection*{Outcomes from this work}\label{outcomes-from-this-work}}
\addcontentsline{toc}{subsubsection}{Outcomes from this work}

\begin{enumerate}
\def\labelenumi{\arabic{enumi}.}
\tightlist
\item
  A clear vision of the challenge you want to tackle.\\
\item
  A set of ideas, papers, and the outcome you ultimately provide.\\
\item
  The evidence you need to work some magic - ie to draw a comic/infographic and then write a very short paper using one of the three formats (narrative, systematic, or meta).
\end{enumerate}

\hypertarget{scicomm}{%
\chapter{Scientific communication}\label{scicomm}}

\includegraphics[width=3in,height=\textheight]{./scicomm.png}

\hypertarget{learning-outcomes-7}{%
\subsection*{Learning outcomes}\label{learning-outcomes-7}}
\addcontentsline{toc}{subsection}{Learning outcomes}

\begin{enumerate}
\def\labelenumi{\arabic{enumi}.}
\tightlist
\item
  Practice your communication skills.
\item
  Be creative.\\
\item
  Simplify complex scientific evidence appropriately.
\end{enumerate}

\hypertarget{context-6}{%
\subsection*{Context}\label{context-6}}
\addcontentsline{toc}{subsection}{Context}

Scientific communication is now its own field of research and activity. It is critical that biology for environmental management promotes better communication. Without effective delivery of evidence, decisions that need science will struggle. Here, we use ten simple rules for \href{https://journals.plos.org/ploscompbiol/article?id=10.1371/journal.pcbi.1005845}{scientific comics} \citep{RN7247} to inform our sharing the scientific synthesis work began in the previous chapter. Simplication of science and aggregation of complex evidence is not a trivial skill. It is a challenge to respectfully communicate science but ensure it is accessible to a wide audience. Comics are perfect. Infographics are also powerful, and they are typically more complex information but compile findings still very visually. A comic tells a story and shares the science whilst infographics are also visual but connect ideas in a mini-poster format to the share evidence. \href{https://www.nature.com/articles/d41586-017-08542-5}{Here are examples} of the very best from the journal Nature.

\hypertarget{instructions-2}{%
\subsection*{Instructions}\label{instructions-2}}
\addcontentsline{toc}{subsection}{Instructions}

This is a stepping stone assignment or scaffolding to recognize the work you have already done in preparation for your short synthesis paper. It also provides an opportunity be more visual, tell a story, be funny, or present the key evidence you found in an infographic.

If you elect to do a comic, you can draw by hand then scan it or sketch using a digital tool.\\
If infographic, try using a free online tool, Keynote on a Mac, or PowerPoint to draw a single-slide and put your key evidence there.

\begin{enumerate}
\def\labelenumi{\arabic{enumi}.}
\tightlist
\item
  Use your evidence from your synthesis.\\
\item
  Pick a few key points. Tell a story visually.
\item
  Save to PDF (only this format will be accepted) and submit to \href{https://www.turnitin.com}{turnitin.com}. Class ID and key provided in the official course outline.
\end{enumerate}

\hypertarget{rubric-1}{%
\subsection*{Rubric}\label{rubric-1}}
\addcontentsline{toc}{subsection}{Rubric}

\begin{tabular}{rllr}
\toprule
item & concept & description & value\\
\midrule
1 & clarity & clear and informative visuals, compelling or funny and accessible graphics & 5\\
2 & science & transparent science such as challenge-to-solution paradigm developed herein & 5\\
3 & synthesis & big ideas put together effectively summarizing the field & 5\\
4 & conclusions & conclusion, punchline, implication stated well & 5\\
11 & total & sum of above & 20\\
\bottomrule
\end{tabular}

\hypertarget{metas}{%
\chapter{Meta-analyses}\label{metas}}

\includegraphics[width=3in,height=\textheight]{./meta.png}

\hypertarget{learning-outcomes-8}{%
\subsubsection*{Learning outcomes}\label{learning-outcomes-8}}
\addcontentsline{toc}{subsubsection}{Learning outcomes}

\begin{enumerate}
\def\labelenumi{\arabic{enumi}.}
\tightlist
\item
  Examine scientific findings from a strength or evidence, big picture viewpoint.
\item
  Practice critical thinkings in evidence reuse in the sciences.\\
\item
  Be able to read evaluate published meta-analyses.
\end{enumerate}

\hypertarget{context-7}{%
\subsubsection*{Context}\label{context-7}}
\addcontentsline{toc}{subsubsection}{Context}

As described in the `scientific synthesis' section in this course, there are at least three contemporary tools available to a synthesis scientist in the environmental sciences. Narrative reviews are most effective when they briefly highlight key findings for a small field of research or identify critical gaps. Systematic reviews are more significant contribution to knowledge engineering because they provide the blueprint for replication, i.e.~they describe how they selected the studies they summarized, and also because they provide more specific evidence. Finally, meta-analyses are the gold standard. For the purposes of this course of study, biology for environmental management, any of these three options are within reach. Nonetheless, metas can be demanding in time, and if you elect to tackle your environmental challenge with this tool, I recommend a very limited and restrictive set of studies. This can be done using very precise terms in The Web of Science or Google Scholar and/or limiting returns of papers to only the top 10-15 most cited or more recent 2-5 years. The capacity to critically examine a published meta is an invaluable asset to becoming a more literate citizen. The paper \href{https://onlinelibrary.wiley.com/doi/10.1002/jrsm.1150}{`How to critically read ecological meta-analyses'}\citep{RN3498} is a resource developed for that purpose. Here is \href{https://onlinelibrary.wiley.com/doi/10.1002/jrsm.1109}{another `how-to-read'} published metas, and finally, if this is the synthesis tool for you, \href{https://link.springer.com/article/10.1007/s11301-021-00247-4}{here is recent practical guide how to do meta-analysis} (in eight steps).

\hypertarget{steps-for-this-component}{%
\subsubsection*{Steps for this component}\label{steps-for-this-component}}
\addcontentsline{toc}{subsubsection}{Steps for this component}

\begin{enumerate}
\def\labelenumi{\arabic{enumi}.}
\tightlist
\item
  Read the suggested how-to paper.\\
\item
  Compare and contrast the evidence you compiled for your synthesis and used for the comic or infographic and that you will use again for the final synthesis paper with the `how to critically evaluate meta' ideas proposed.\\
\item
  Decide what is most appropriate based on the evidence and reporting provided in those peer-reviewed studies.\\
\item
  Skip ahead to the `paper writing' section of this course and the `paper' section too for rubric and instructions to facilaite the best decision for what category of paper you select to write. All the very, very short, but each differs slightly in they handle and present evidence.
\end{enumerate}

\hypertarget{writing}{%
\chapter{Paper writing}\label{writing}}

\includegraphics[width=3in,height=\textheight]{./paperwriting.png}

\hypertarget{learning-outcomes-9}{%
\subsubsection*{Learning outcomes}\label{learning-outcomes-9}}
\addcontentsline{toc}{subsubsection}{Learning outcomes}

\begin{enumerate}
\def\labelenumi{\arabic{enumi}.}
\tightlist
\item
  Critically weigh evidence to decide how to best report it.
\item
  Write a short synthesis paper.\\
\item
  Consolidate and hone your writing to be logical and structured even when short.
\end{enumerate}

\hypertarget{context-8}{%
\subsubsection*{Context}\label{context-8}}
\addcontentsline{toc}{subsubsection}{Context}

Writing is an art. Telling a story well, even in science, is a fundamental professional life skill. Being brief is a virtue. Integrating evidence into scientific writing without being overwhelming is ideal. Here, you can tackle this by reusing your environmental evidence for the challenge you selected. This is an excellent, brief set of guidelines to consider, in general, to consider for better, more clearly structured writing \citep{RN7239}. Yes, there is always room for improvement.

Specific this course, you must choose narrative, systematic, or meta. This might have been decided right-up-front because of the nature of the challenge and the evidence that was reported in the papers you compiled. For instance, there was limited data you could extract from the papers, each paper tackled the same challenge but did so very differently, or it was not tractable to process the evidence more formally either then identifying what we know, do not, or where to go.

\hypertarget{steps-for-this-component-1}{%
\subsubsection*{Steps for this component}\label{steps-for-this-component-1}}
\addcontentsline{toc}{subsubsection}{Steps for this component}

\begin{enumerate}
\def\labelenumi{\arabic{enumi}.}
\tightlist
\item
  Find your style.\\
\item
  Organize your thoughts.
\item
  Decide if you need a single, simple graphic to support your paper.\\
\item
  Write your draft, rapidly, now. Correct later. Let is sit for while. Aim for \textless2500 words.\\
\item
  See final instructions and rubric in next section.
\end{enumerate}

\hypertarget{paper}{%
\chapter{Paper}\label{paper}}

\includegraphics[width=3in,height=\textheight]{./paper.png}

\hypertarget{learning-outcomes-10}{%
\subsection*{Learning outcomes}\label{learning-outcomes-10}}
\addcontentsline{toc}{subsection}{Learning outcomes}

\begin{enumerate}
\def\labelenumi{\arabic{enumi}.}
\tightlist
\item
  Produce a clear, succinct summary of research.\\
\item
  Demonstrate mastery of evidence implementation for the environmental sciences.\\
\item
  Appreciate the challenge of effective simplification in the sciences.
\end{enumerate}

\hypertarget{choices-redux}{%
\subsection*{Choices redux}\label{choices-redux}}
\addcontentsline{toc}{subsection}{Choices redux}

\begin{enumerate}
\def\labelenumi{(\arabic{enumi})}
\item
  Narrative review as a highlight, short commentary, or new idea paper that is a snapshot of the key findings from a field research summarizing the main discoveries and/or listing the most critical research gaps. Papers like these are often called \textbf{Insights}, \textbf{Forum}, or \textbf{Ignite}.
\item
  Systematic reviews are similar to narrative reviews, but clear criteria are listed explaining how you selected papers (i.e.~these search terms used in the Web of Science, only studies that had these key inclusion attributes, etc.). Systematic are more replicable because others could follow your steps and get the same set of papers and hopefully reach similar conclusions about the corpus of evidence. These reviews also typically provide some simple quantitative data about the set of studies such as number of countries where the research was done, total sample sizes, number of variables examined, or any attributes that describe what the research was for a specific detail. The narrative component is similar to the first option because it can state what we know and what do not but these reviews do so \emph{much} more precisely. Even a few numbers go a long way to convincing people about the extent that we know or have studied a subject in science. Papers like these are often termed \textbf{Short Commentary} or \textbf{Mini-Review}.
\item
  Meta-analyses are systematic reviews plus for each primary study you summarize, you capture the relative efficacy of the treatment tested. Papers like this are often termed \textbf{Reviews} or \textbf{Meta-analyses} but other terms can be used too. Note: in some fields of research the terms `systematic review' and `meta-analysis' are used interchangeably, but in most environmental sciences, meta-analyses always have a measure of the strength of evidence from each studied included in the synthesis whilst systematic reviews typically do not.
\end{enumerate}

\begin{tabular}{llll}
\toprule
synthesis & elements & benefits & limitations\\
\midrule
narrative review & summary, insights, next steps & more opinion, can be shorter, less detailed processing the research literature & can be less compelling without some specific evidence listed and is difficult to repeat\\
systematic review & summary, insights, next steps, explanation of how studies were selected in the review, can have counts of ideas tested & more specific, can be repeated, and can be more compelling & need to keep track of how you selected papers, need to sort through papers in more detail in addition to capturing main ideas\\
meta-analyses & all of above possible but must also include an assessment of the strength of evidence of each study included in the synthesis & gold standard, reader can get a sense of how effective a treatment or intervention is relative to another & need to extract means or measures of efficacy with sample sizes or variance from each study\\
\bottomrule
\end{tabular}

\hypertarget{bottom-line}{%
\subsection*{Bottom line}\label{bottom-line}}
\addcontentsline{toc}{subsection}{Bottom line}

\textbf{Narrative review} = Insights, Forum, or Ignite style paper common in current journals.\\
Max 2000 words (not including lit cited)\\
10-15 references

Highlight what we know, what we do not or a key gap we need to study (ie insights)

\textbf{Systematic review} = Short Commentary or Mini-review format in journals.\\
Max 2500 words (not including lit cited)\\
10 references

A few key quantitative details of the research literature such the frequency a specific concept is studied, total sample sizes or counts of variables measured. This is a comment on what we know with some evidence in numbers.

\textbf{Meta-analysis} = Review or Meta-analysis format paper.\\
2500 words (not including lit cited)\\
One figure plotting main finding of mean effects, i.e.~a \href{https://en.wikipedia.org/wiki/Forest_plot}{``Forest plot''}\\
15 references

Find 5-10 papers that have the same data reported in each that you can use, get the means, r2, n or whatever is reported similarly, calculate a simple \href{https://en.wikipedia.org/wiki/Effect_size}{effect size} such as ratio of treatment over control or a proportion that an intervention helped, plot this outcome, and write short paper commenting how well an intervention in the environmental sciences works (i.e.~adding water to restore deserts or excluding cattle, one direct and simple variable that you care about)

\textbf{Whatever format you elect to use} the goal is to provide a super short \& sweet synthesis of the research and clearly demonstrate that you used scientific evidence from papers that critically informs better environmental management choices.

\hypertarget{instructions-3}{%
\subsection*{Instructions}\label{instructions-3}}
\addcontentsline{toc}{subsection}{Instructions}

\begin{enumerate}
\def\labelenumi{\arabic{enumi}.}
\item
  Select your format for paper. Here is an excellent example of \href{https://www.oikosjournal.org/authors/author-guidelines}{guidelines} for a world-class journal. ALL are short.
\item
  Find a few examples of that type of paper. Read them.
\item
  Include your name, surname, and student ID at the top of each page.
\item
  Strict adhere to word limits. A total of three short paragraphs, no more.
\item
  Write your paper in preferred text editor, print/save to PDF, save and submit to \href{https://www.turnitin.com}{turnitin.com}. Class ID and key provided in the official course outline.
\end{enumerate}

\hypertarget{rubric-2}{%
\subsection*{Rubric}\label{rubric-2}}
\addcontentsline{toc}{subsection}{Rubric}

\begin{tabular}{rllr}
\toprule
item & concept & description & value\\
\midrule
1 & title & clear and relevant & 2\\
2 & introduction & one paragraph, short and direct, states the environmental challenge, why we should care, how to potentially solve it & 8\\
3 & main text & one paragraph, state main findings for whatever format you selected & 10\\
4 & implications & final \& single paragraph as well, state the meaning of what you found, how it can used, who will benefit, how to use knowledge or gaps thereof & 10\\
5 & literature cited & 10-15 references depending on type of paper, relevant, useful evidence, mostly recent i.e. last 5 years, clear choices & 5\\
\addlinespace
6 & transparency & clear writing, results, counts, plot or data clear if included, transparency in how evidence was used, balanced yet critical thinking review and synthesis & 5\\
11 & total & sum of above & 40\\
\bottomrule
\end{tabular}

\hypertarget{create}{%
\chapter{CREATE}\label{create}}

\includegraphics[width=3in,height=\textheight]{./castle.png}

\hypertarget{learning-outcomes-11}{%
\subsection*{Learning outcomes}\label{learning-outcomes-11}}
\addcontentsline{toc}{subsection}{Learning outcomes}

\begin{enumerate}
\def\labelenumi{\arabic{enumi}.}
\tightlist
\item
  Creatively read scientific literature.\\
\item
  Hone your critical thinking.\\
\item
  Develop a new set of tools for processing readings.
\end{enumerate}

\hypertarget{create-philosophy}{%
\subsection*{CREATE philosophy}\label{create-philosophy}}
\addcontentsline{toc}{subsection}{CREATE philosophy}

Science is creative. Reading scientific literature can also be a creative process. \href{https://help.open.ac.uk/active-reading}{Active reading} includes highlights, annotation, and note taking. These approaches to reading and study can increase retention, comprehension, and engagement with scientific ideas. However, the teaching philosophy and pedagogy of CREATE is novel and relevant and takes this one step further. \href{https://academic.oup.com/genetics/article/176/3/1381/6062159}{\textbf{CREATE pedagogy}} stands for \textbf{C}onsider, \textbf{R}ead, \textbf{E}lucidate the hypothesis or purpose, \textbf{A}nalyze and interpret data or evidence, and \textbf{T}hink of the next \textbf{E}xperiment. This reading process is creative, active, and phenomenal practice for all professional work. If this process seems useful to you, please give it a try for this course. This work is \textbf{optional}, but it is literally only 5-10 minutes extra per week. Here are some proposed exercises for you to consider to test out the CREATE approach. If there is time in lectures, we will do them together. Use paper, powerpoint, or any sketching digital tool you prefer to write or draw. If you elect to try them out and want to submit them as well, \textbf{optional}, there are spots provided in \href{https://turnitin.com}{turnitin.com}.

\hypertarget{create-applications}{%
\subsection*{CREATE applications}\label{create-applications}}
\addcontentsline{toc}{subsection}{CREATE applications}

\begin{tabular}{rll}
\toprule
week\_ID & CREATE\_app & description\\
\midrule
1 & new title \& abstract & write a catchy, compelling alternative title to the reading and a very short 250 word abstract in your own words\\
2 & concept map & draw a Venn diagram, flowchart, or any type of visual that summarizes how the big ideas from the reading this week connect to one another\\
3 & cartoon & draw a fun cartoon summarizing the reading this week\\
4 & questions & a good reading or paper should generate as many new questions as the ones it answers, list a few for the reading this week\\
5 & made-up data & sketch a plot of data you would like to see supporting the main idea or hypothesis in the reading this week\\
\addlinespace
6 & experiment & sketch a next step experiment of any sort to support the reading this week\\
\bottomrule
\end{tabular}

\hypertarget{life}{%
\chapter{Life}\label{life}}

\includegraphics[width=3in,height=\textheight]{./life.png}

\hypertarget{action}{%
\subsection*{Action}\label{action}}
\addcontentsline{toc}{subsection}{Action}

Make it count! Help those around you use scientific evidence to inform their decisions. Science can support better decisions in health, the environment, social policy, and all domains of the human endeavour. We need to make better decisions, and we can do it. Use your creative spirit and now finely-honed critical thinking. Be on the lookout for solutions. Frame problems in life, big or small, as challenges. Seek solutions in the obvious and in the novel. Less is more. Often the simplest fix is the best. Good luck, and I hope you enjoyed these materials and this course as much I did.

  \bibliography{book.bib,packages.bib}

\end{document}
